\section{AI Usage in Project Development}

\subsection{Marty Lauterbach}

In my workflow to create the architectural design, I started off by reading up on documentation of react native and expo online. When I established my first pattern, I used claudes web interface to analyse my planning with his deep research tool, to refine my design with its suggested sources. This helped me widen my horizon of avaliable resources and let me judge the best pattern for our requirements and the three main parts of it.

In my research to find all possible ways to implement :

A. holding the logging in context while the user captures media

B. react native camera, audio and video recording

C. Machine learning models for object, image and audio recognition

D. Localization

E. OpenStreetMap Integration

F. Complete Bird .CSV Data -> SQL Lite

G. Wikimedia latin name based picture scraping

H. Dynamic bundling of assets to the app

-> I used a cycle of personal research, planning, redefining based on deeper research by Claude, prototyping hand-in-hand and building testing suites with python and javascript with Claude Code (local API Agent) on Manjaro Arch Linux, testing implementation, improving UX Design after testing, running testing suites to ensure full integration of global theme and components, compressing the created architecture, research and design into Markdown-files (by Claude) in /dev for team documentation.

This helped me to plan more throughly and to create a more robust scaleable modularized architecture, while maintaining thorough tests that adhered to my established global requirements like theme, user stories and design principles.

Furthermore, I used Claude Code to create research skimming through locally cloned public repositories of machine learning in kotlin (whoBird), to recognize and tell me the architectural layout, the input tensors and the output label structure of the .tflite models that are used for the audio pipeline. This let me understand the needed react native structure for using the models ourself.
I then created and extended a testing suite to test different inputs into the models to find the best possible filtering methods and audio decodings with mel spectograms to input into the model.

For the laTex Documentation, I let Claude Code help me in creating python scripts to compare the documentation to the actual app to find unreferenced technical claims that would need file and line references, as well as possible pseudo-code additions.

In Conclusion, AI was not good enough to be more than a helpful tool for this project (in research and troubleshooting with interfaces) since it did not grasp the ML pipelines, the correct usage of context, the UX Design, the UI constraints and the always evolving, and therefore being outside its mostly 2024 scope, landscape of react native tensorflow (that has since for example been deprecated in Hugging Face).

Thanks,
Marty

\subsection{Youmna Samouneh}
I used AI to implement the translations.

\subsection{Luis Wehrberger}

